\documentclass[11pt]{article}
\usepackage[margin=1in]{geometry}
\usepackage{amsmath,amssymb}
\usepackage{tikz-feynman}
\usepackage{graphicx}
\usepackage{hyperref}
\usepackage{slashed}
\usepackage{booktabs}

\newcommand{\slashedp}[1]{\slashed{#1}}

\title{Reproducing Vermaseren's numerical-instability claim for\\
$\gamma^*\gamma^*\to\mu^+\mu^-$}
\author{nikhef-form-lab, Part 5 (\texttt{06-mc-multiperipheral})}
\date{\today}

\begin{document}
\maketitle

\begin{abstract}
J.A.M.\ Vermaseren's \emph{Two Photon Processes at Very High Energies}
(NIKHEF-H/82-15, 1982) states, discussing the reaction
$e^+e^-\to e^+e^-\mu^+\mu^-$ at PETRA/LEP energies: ``If the matrix element
were to be completely evaluated in the standard fashion -- i.e.\ expressing
it completely in terms of 4-vector products and then substituting their
numerical values for each Monte Carlo generated point in phase space --
then the cancellations between the various terms would be so bad that even
the 60 bit accuracy of a CDC computer would not suffice.'' This report
reproduces that claim end to end for the two-photon subprocess
$\gamma^*\gamma^*\to\mu^+\mu^-$: deriving the squared matrix element by hand
via the naive (``standard fashion'') route and via Vermaseren's
Schouten-identity-stabilized route, cross-checking both derivations
symbolically, and then evaluating both forms numerically in IEEE double
precision at kinematic points approaching the multiperipheral peak
($t_1,t_2\to0$) to measure directly how many significant digits are lost
in the naive route.
\end{abstract}

\tableofcontents

\section{The process and the two diagrams}
\label{sec:diagrams}

We study the subprocess $\gamma^*(q_1)\,\gamma^*(q_2)\to\mu^-(p_6)\,\mu^+(p_7)$,
following Vermaseren's Fig.~3 labeling (reproduced in
Figures~\ref{fig:diagram-k1}--\ref{fig:diagram-k2}). This is the
matrix-element core of the full $e^+e^-\to e^+e^-\mu^+\mu^-$ two-photon
process (Fig.~1 of the same paper): the electron and positron each radiate
a virtual photon ($q_1$, $q_2$), which then fuse via muon exchange into the
$\mu^+\mu^-$ pair. The electron/positron lines and the two photon
propagators $1/q_1^2,\,1/q_2^2$ are common to both diagrams and are not the
source of the numerical instability discussed here; the instability lives
entirely in the $\gamma^*\gamma^*\to\mu^+\mu^-$ subprocess itself, so we
isolate it.

There are exactly two tree diagrams, differing by which photon attaches to
the muon line first:

\begin{itemize}
  \item \textbf{Diagram $k_1$}: photon $q_1$ attaches first; the internal
    muon propagator carries momentum $k_1 = p_6-q_1 = q_2-p_7$.
  \item \textbf{Diagram $k_2$} (crossed): photon $q_2$ attaches first;
    the internal muon propagator carries momentum $k_2 = p_6-q_2 = q_1-p_7$.
\end{itemize}

\begin{figure}[htbp]
  \centering
  \includegraphics[width=0.45\textwidth]{figures/mumu_diagram_k1}
  \caption{Diagram $k_1$: $q_1$ attaches to the muon line first.}
  \label{fig:diagram-k1}
\end{figure}

\begin{figure}[htbp]
  \centering
  \includegraphics[width=0.45\textwidth]{figures/mumu_diagram_k2}
  \caption{Diagram $k_2$ (crossed): $q_2$ attaches to the muon line first.}
  \label{fig:diagram-k2}
\end{figure}

\subsection{Kinematic conventions}

On-shell conditions: $p_6^2=p_7^2=m^2$ (muon mass), while the photons are
virtual with $q_1^2=t_1$, $q_2^2=t_2$ (generally $t_1,t_2\neq0$, since these
are internal propagators of the full process, not real external photons).
Momentum conservation at the $\gamma^*\gamma^*$ vertex reads
$q_1+q_2=p_6+p_7\equiv Q$. We will also use $k_1^2=(p_6-q_1)^2$ and
$k_2^2=(p_6-q_2)^2$, with propagator denominators
$D_1\equiv k_1^2-m^2$, $D_2\equiv k_2^2-m^2$.

The ``multiperipheral peak'' regime the paper's warning is about is
$t_1,t_2\to0$ (both photons nearly on-shell/real) at fixed, large $Q^2=s$
-- the same corner of phase space studied numerically for the $\pi^0$ case
in Topic 5.

\section{Feynman rules and the amplitude}
\label{sec:feynman-rules}

Using standard QED Feynman rules (photon-fermion vertex $-ie\gamma^\mu$,
fermion propagator $i(\slashed k+m)/(k^2-m^2)$):

\begin{align}
  \mathcal M_1 &= (-ie)^2\,\bar u(p_6)\,(-i\slashed e_1)\,
    \frac{i(\slashed k_1+m)}{D_1}\,(-i\slashed e_2)\,v(p_7)
    \label{eq:M1} \\
  \mathcal M_2 &= (-ie)^2\,\bar u(p_6)\,(-i\slashed e_2)\,
    \frac{i(\slashed k_2+m)}{D_2}\,(-i\slashed e_1)\,v(p_7)
    \label{eq:M2}
\end{align}

Dropping the overall coupling and bookkeeping factors of $i$ into a common
prefactor (matching Vermaseren's own convention, cf.\ formula (III.3) of
\texttt{twophoton.pdf}), the total amplitude is

\begin{equation}
  \mathcal M \;\propto\; \bar u(p_6)\,N\,v(p_7), \qquad
  N \equiv \slashed e_1\,\frac{\slashed k_1+m}{D_1}\,\slashed e_2
    \;+\; \slashed e_2\,\frac{\slashed k_2+m}{D_2}\,\slashed e_1
  \label{eq:N-def}
\end{equation}

This matches the left-hand side of formula (III.3) in \texttt{twophoton.pdf}
exactly, confirming the starting point before proceeding to the trace
algebra.

\section{The naive (``standard fashion'') squared amplitude}
\label{sec:naive}

\emph{[To be filled in: full spin-summed trace
  $\mathrm{Tr}[(\slashed p_6+m)\,N\,(\slashed p_7-m)\,\bar N]$, expanded
  into its 16 cross terms, each reduced via the standard trace identities.
  This is the ``expressing it completely in terms of 4-vector products''
  route the paper's quoted paragraph describes.]}

\section{Vermaseren's stable form}
\label{sec:stable}

\emph{[To be filled in: the Schouten-identity rewrite reproducing formula
  (III.5)/(III.10)--(III.11), independently re-derived rather than quoted.]}

\section{Symbolic cross-check}
\label{sec:symbolic-check}

\emph{[To be filled in: sympy verification that Sections~\ref{sec:naive}
  and \ref{sec:stable} are algebraically identical at generic symbolic
  kinematics.]}

\section{Numerical demonstration of the cancellation}
\label{sec:numerics}

\emph{[To be filled in: double-precision evaluation of both forms at
  kinematic points approaching $t_1,t_2\to0$, with explicit digit-loss
  counts, following the same style as
  \texttt{kinc2-driver/check\_gram\_minors.c} for the $\pi^0$ case.]}

\section{Monte Carlo scan}
\label{sec:mc}

\emph{[To be filled in: fraction of randomly sampled phase-space points at
  which the naive route loses catastrophic precision, as a function of
  proximity to the multiperipheral peak.]}

\section{Conclusion}
\label{sec:conclusion}

\emph{[To be filled in once Sections~\ref{sec:naive}--\ref{sec:mc} are
  complete.]}

\end{document}
